\begin{figure}
    \begin{lstlisting}[]
void spmspm(row_ptrs_A: mref<? x idx>,
            col_idxs_A: mref<? x idx>,
            vals_A: mref<? x f32>,
            row_ptrs_B: mref<? x idx>,
            col_idxs_B: mref<? x idx>,
            vals_B: mref<? x f32>,
            out: CSR) {
  for (idx i=0; i<n_rows; i++) {
    clear(touched);
    for (idx p=row_ptrs_A[i]; p<row_ptrs_A[i+1]; p++) {
      idx k = col_idxs_A[p];
      f32 a = vals_A[p];
      for (idx q=row_ptrs_B[k]; q<row_ptrs_B[k+1]; q++) {
        idx j = col_idxs_B[q];
        if (!mark[j]) {
          mark[j] = true;
          touched.push(j);
          accum[j] = 0;
        }
        accum[j] += a * vals_B[q];
      }
    }
    sort(touched); // Optional: keep CSR columns sorted
    for (idx j : touched) {
      emit(out, i, j, accum[j]);
      mark[j] = false;
    }
  }
}
    \end{lstlisting}
    \vspace{-1em}
    \caption[Pseudocode of sparse matrix-sparse matrix multiplication.]{Pseudocode for Gustavson-style SpMSpM. The outer loop builds one output row at a time. For each nonzero $A_{ik}$, the middle loop uses $k$ to look up row $k$ of $B$, and the inner loop scans that row. Products are accumulated in a row-local sparse accumulator: \texttt{mark} identifies whether a column has already been touched, and \texttt{touched} records the nonzero columns of the output row. The output row can therefore be emitted by iterating only over \texttt{touched}, avoiding a scan of the full dense accumulator. Compared with SpMV, SpMSpM performs scan-and-lookup over whole sparse rows and introduces accumulation of intermediate results.}
    \label{fig:background-spmspm-pseudocode}
\end{figure}
